Parkinson's disease (PD) affects over 10 million individuals worldwide\cite{dorsey2018global}, with substantial heterogeneity in progression trajectories necessitating personalized prognostic models\cite{fereshtehnejad2017subtypes}. Graph neural networks (GNNs) have emerged as powerful tools for integrating multimodal patient data—clinical assessments, neuroimaging biomarkers, and genetic variants—through learned patient similarity networks\cite{zhou2020graph,parisot2018disease}. However, the ``black box'' nature of deep learning models remains a critical barrier to clinical adoption\cite{amann2020explainability,holzinger2019causability}, particularly for high-stakes medical decisions where transparency and interpretability are paramount\cite{rudin2019stop}.

Explainable AI (XAI) methods for GNNs are nascent compared to traditional neural networks\cite{yuan2021explainability}. While attention mechanisms\cite{veličković2018graph} provide some interpretability, comprehensive frameworks validating GNN predictions against clinical knowledge are lacking. Existing medical AI explainability often relies on single methods—saliency maps\cite{selvaraju2017grad} or feature importance scores\cite{lundberg2017unified}—risking method-specific artifacts. Multi-method validation, analogous to triangulation in qualitative research, strengthens evidentiary claims by demonstrating convergent findings across complementary analytical approaches\cite{molnar2020interpretable}.

We address this gap with a comprehensive explainability framework for the Graph-Integrated Multimodal Attention Network (\giman)\cite{authorYYYY_giman}, a graph attention network (GAT) trained on 2,046 Parkinson's Progression Markers Initiative (PPMI) participants. Our framework integrates six complementary methods applied to three prediction tasks: (1) diagnostic classification (healthy control vs. PD vs. SWEDD), (2) progression subtype stratification (fast/moderate/slow progressors), and (3) prodromal-to-clinical conversion prediction. The methods span structural (attention weights, GNNExplainer subgraphs), attribution-based (IntegratedGradients, GradientSHAP), embedding-based (patient clustering), and counterfactual approaches, providing orthogonal perspectives on model behavior.

Our contributions are threefold. \textbf{First}, we demonstrate that multi-method explainability reveals convergent evidence for clinically meaningful predictors, with motor progression rate (UPDRS slope) consistently ranked highest across tasks. \textbf{Second}, we show that GNN attention mechanisms learn interpretable patient similarity (88\% same-label coherence), validating the patient similarity graph paradigm. \textbf{Third}, we provide actionable, patient-level explanations through counterfactual analysis and clinical dashboards, bridging the gap between model predictions and clinical decision-making. This framework establishes a roadmap for trustworthy GNN deployment in precision medicine.
